\subsection{Application: the Laplace transformation}

In this No., we denote by M a commutative monoid, whose law of composition is written additively, equipped with the topology of a locally compact space, for which the mapping \((m, m') \mapsto m + m'\) of \(M \times M\) into M is continuous. The neutral element of M is denoted by 0. One calls character of M every bounded continuous complex function \(\chi\) on M satisfying the relations

\[
\chi (m + m ^ {\prime}) = \chi (m) \cdot \chi (m ^ {\prime}), \quad \chi (0) = 1, \quad | \chi (m) | \leqslant 1\tag{7}
\]

for \(m, m'\) in M. If \(\chi\) and \(\chi'\) are characters, then so is \(\chi\chi'\) . The set of characters of M is a monoid, denoted X; equip it with the topology of compact convergence, for which the mapping \((\chi, \chi') \mapsto \chi\chi'\) of X × X into X is continuous. The neutral element of X is the constant function 1.

For every bounded complex measure \(\mu\) on M, one calls Laplace transform of \(\mu\) the function \(\mathcal{L}\mu\) on X defined by

\[
(\mathscr {L} \mu) (\chi) = \int_ {\mathrm{M}} \chi (m) d \mu (m).\tag{8}
\]

By Th. 3 of GT, X, §3, No.~4, the mapping \((m,\chi)\mapsto\chi(m)\) of M×X into C is continuous and bounded. The corollary of Prop. 13 of No.~6 then implies the following result:

PROPOSITION 14. --- For every bounded complex measure \(\mu\) on M, the function \(\mathcal{L}\mu\) on X is continuous and bounded. If \(\mathcal{M}_{+}^{b}(M)\) is equipped with the tight topology and \(\mathcal{C}^{b}(X; \mathbf{C})\) with the topology of compact convergence, the mapping \(\mu \mapsto \mathcal{L}\mu\) of \(\mathcal{M}_{+}^{b}(M)\) into \(\mathcal{C}^{b}(X; \mathbf{C})\) is continuous.

The set of characters of M that tend to 0 at infinity will be denoted \(X_{0}\) ; this set is stable under multiplication. We shall say that a submonoid \(^{(1)}\) S of X is full if S is stable for the mapping \(\chi\mapsto\overline{\chi}\) , \(S\cap X_{0}\) separates the points of M (GT, X, §4, No.~1, Def. 1) and if, given any \(m\in M\) , there exists an element \(\chi\) of \(S\cap X_{0}\) such that \(\chi(m)\neq0\) .

Suppose in addition that M is a noncompact abelian group. Let f be a function on M that tends to 0 at infinity; the same is then true of the function \(x \mapsto f(x)f(-x)\) on M, whereas every character \(\chi\) of M satisfies \(\chi(x)\chi(-x) = \chi(0) = 1\) . It follows that \(X_{0}\) is empty, and that X does not contain any full submonoid. Thus, Theorem 3 below does not apply to locally compact groups that are not compact.

THEOREM 3. --- Let S be a full submonoid of X.

\begin{enumerate}
\def\labelenumi{\alph{enumi})}
\item
  If \(\mu\) and \(\mu'\) are two bounded complex measures on \(\mathbf{M}\) , such that \(\mathcal{L}\mu\) and \(\mathcal{L}\mu'\) have the same restriction to \(\mathrm{S} \cap \mathrm{X}_0\) , then \(\mu = \mu'\) .
\item
  Let \(\mathfrak{F}\) be a filter on \(\mathcal{M}_{+}^{b}(\mathbf{M})\) , such that \(\mathcal{L}\lambda(s)\) has a limit \(\Phi(s) \in \mathbf{C}\) with respect to \(\mathfrak{F}\) for every \(s \in \mathbf{S}\) . Then the filter \(\mathfrak{F}\) converges vaguely to a bounded positive measure \(\mu\) , and \(\Phi(s) = \mathcal{L}\mu(s)\) for all \(s \in \mathbf{S} \cap \mathbf{X}_0\) .
\item
  Under the hypotheses of \(b\) ), suppose in addition that the closure of \(S \cap X_0\) contains 1, and that the function \(\Phi\) on \(S\) is continuous at the point 1. Then \(\mathfrak{F}\) converges tightly to \(\mu\) , and \(\Phi(s) = \mathcal{L}\mu(s)\) for all \(s \in S\) .
\end{enumerate}

We shall denote by E the algebra of continuous complex functions tending to 0 at infinity on M, and by A the linear subspace of E generated by \(S \cap X_0\) ; then \(\mathfrak{A}\) is a subalgebra of \(E\) stable under the mapping \(f \mapsto \overline{f}\) ; since \(S\) is a full submonoid of \(X\) , Cor. 2 of Prop. 7 of GT, \(X\) , §4, No.~4 implies that \(\mathfrak{A}\) is dense in \(E\) .

Let us prove \(a\) : by hypothesis, \(\mu(f) = \mu'(f)\) for every \(f \in \mathfrak{A}\) ; since \(\mu\) and \(\mu'\) are continuous linear forms on \(\mathbf{E}\) , this implies that \(\mu(f) = \mu'(f)\) for \(f \in \mathbf{E}\) , and in particular for every continuous function \(f\) with compact support, whence \(\mu = \mu'\) .

Let us place ourselves under the hypotheses of \(b\) ). The number \(\Phi(1) = \lim_{\lambda, \mathfrak{F}} \lambda(1)\) is real and positive; let there be given a real number \(a > \Phi(1)\) ; since

\(\| \lambda \| = \mathcal{L}\lambda (1)\) for \(\lambda \in \mathcal{M}_{+}^{b}(\mathrm{M})\) , the relation \(\lim_{\lambda ,\mathfrak{F}}\mathcal{L}\lambda (1) = \Phi (1)\) implies that the set H of measures \(\lambda \in \mathcal{M}_{+}^{b}(\mathrm{M})\) such that \(\| \lambda \| \leqslant a\) belongs to \(\mathfrak{F}\) . Since \(\mathcal{M}^b (\mathrm{M};\mathbf{C})\) may be identified with the dual of the normed space E (Ch. III, §1, No.~8 \& §1, No.~2, Prop. 3), the space H is compact for the topology \(\sigma (\mathcal{M}^b (\mathrm{M};\mathbf{C}),\mathrm{E})\) . On the other hand (TVS, III, §3, No.~4, Prop. 5), this topology coincides on H with the topology of pointwise convergence in any total subset of E. In particular, since \(\mathfrak{A}\) is dense in E, and the same is true of the space of continuous functions with compact support (Ch. III, §1, No.~2, Prop. 3), the topology of pointwise convergence in \(\mathrm{S}\cap \mathrm{X}_0\) coincides on H with the vague topology, and H is compact for this topology. It follows at once that \(\mathfrak{F}\) converges vaguely to a measure \(\mu \in \mathrm{H}\) , and that \(\mathcal{L}\mu (s) = \lim_{\lambda ,\mathfrak{F}}\mathcal{L}\lambda (s)\) for all \(s\in \mathrm{S}\cap \mathrm{X}_0\) .

Finally, let us pass to \(c\) . Since the functions \(\Phi\) and \(\mathcal{L}\mu\) are continuous at the point \(1 \in S\) , and equal on \(S \cap X_0\) , and since 1 is in the closure of \(S \cap X_0\) , we have \(\Phi(1) = \mathcal{L}\mu(1)\) . In other words, \(\lim_{\lambda, \mathfrak{F}} \lambda(1) = \mu(1)\) . Prop. 9 of No.~3 then shows that \(\mu\) is the tight limit of the filter \(\mathfrak{F}\) . Every element of \(S\) being a bounded continuous function on \(M\) , this implies that \(\Phi(s) = \lim_{\lambda, \mathfrak{F}} \lambda(s) = \mu(s) = \mathcal{L}\mu(s)\) for all \(s \in S\) .

COROLLARY. --- Let S be a full submonoid of X, such that the closure of \(S \cap X_{0}\) contains 1. Let L be the subset of \(\mathcal{C}^{b}(S; \mathbf{C})\) consisting of the restrictions to S of the Laplace transforms of the measures \(\lambda \in \mathcal{M}_{+}^{b}(M)\) .

\begin{enumerate}
\def\labelenumi{\alph{enumi})}
\item
  The set L is closed in the space \(\mathcal{C}^{b}(\mathrm{S};\mathbf{C})\) equipped with the topology of pointwise convergence.
\item
  The mapping \(\lambda \mapsto (\mathcal{L}\lambda)_\mathrm{S}\) is a homeomorphism of \(\mathcal{M}_+^b (\mathbf{M})\) onto \(\mathbf{L}\) , if \(\mathcal{M}_+^b (\mathbf{M})\) is equipped with the tight topology and \(\mathbf{L}\) with the topology of pointwise convergence.
\item
  The topology of pointwise convergence and the topology of compact convergence coincide on L.
\end{enumerate}

The assertions \(a)\) and \(b)\) are immediate consequences of Th. 3; the assertion \(c)\) follows from \(b)\) and Prop. 14, since the topology of compact convergence is finer than that of pointwise convergence.

One must be on guard that L is not closed in the set of all bounded complex functions on S, equipped with the topology of pointwise convergence. Assume for example the notations of Example 2 below (M = R+, S identified with R+). The Laplace transforms of the measures \(\varepsilon_{n}\) ( \(n \in N\) ) are the functions \(t \mapsto e^{-nt}\) on R+; as n tends to \(+\infty\) , these functions converge pointwise to the function equal to 1 for t = 0 and to 0 for \(t \neq 0\) , which does not belong to L.

Example 1). --- Take for M the set N of positive integers, equipped with the law of addition and with the discrete topology. Let D be the unit disc of C (the set of complex numbers of absolute value ≤ 1) equipped with the topology induced by C and with the law induced by multiplication. For every z ∈ D, let us denote by f(z) the character n ↦ z\^{}n of N. For every character χ of N, denote by g(χ) the complex number χ(1) ∈ D. One verifies immediately that f and g are mutually inverse homeomorphisms between D and X, which will permit us, from now on, to identify X and D. The set of characters tending to 0 at infinity may then be identified with the set D₀ of complex numbers of absolute value \textless{} 1. Finally, the interval {[}0, 1{]} of R is a full submonoid of D, and 1 is in the closure of {]}0, 1{]} ∩ D₀ = {]}0, 1{[}.

Every measure \(\mu\) on \(\mathbf{N}\) may be written in a unique way in the form \(\mu = \sum_{n\in \mathbf{N}}u_n\cdot \varepsilon^n\) , and \(\mu\) is bounded if and only if \(\sum_{n}|u_n| < +\infty\) ; one then has \(\mathcal{L}\mu (z) = \sum_{n\in \mathbf{N}}u_nz^n\) for \(z\in \mathbf{D}\) . This function is continuous on \(\mathbf{D}\) ; it is customary to call it the generating function of the summable sequence \((u_n)_{n\in \mathbf{N}}\) . Transcribed into this language, Th. 3 yields the following result (taking into account Prop. 9 of No.~3):

PROPOSITION 15. --- Let A be a set equipped with a filter F. For every \(\alpha\in A\) , let \((u_{\alpha,n})_{n\in\mathbf{N}}\) be a summable sequence of positive numbers, and let \(\Phi_{\alpha}\) be the function defined on the interval {[}0,1{]} of R by \(\Phi_{\alpha}(x)=\sum_{n\in\mathbf{N}}u_{\alpha,n}x^{n}\) . In order that there exist a summable sequence \((u_{n})_{n\in\mathbf{N}}\) of positive numbers such that

\[
\lim _ {\alpha , \mathfrak {F}} u _ {\alpha , n} = u _ {n} \text {   for   all   } n, \quad \lim _ {\alpha , \mathfrak {F}} \sum_ {n \in \mathbf {N}} u _ {\alpha , n} = \sum_ {n \in \mathbf {N}} u _ {n},
\]

it is necessary and sufficient that \(\Phi_{\alpha}\) converge pointwise on \(]0,1]\) , with respect to \(\mathfrak{F}\) , to a function \(\Phi\) continuous at the point 1. In this case, \(\Phi(x) = \sum_{n \in \mathbf{N}} u_n x^n\) for all \(x \in ]0,1]\) .

Analogous results are obtained by taking M to be the monoid \(N^{n}\) , where n denotes an integer \textgreater1; the space X may then be identified with \(D^{n}\) , and one can choose \([0,1]^{n}\) as the full submonoid. We leave to the reader the task of transcribing Th. 3 in this case.

Example 2). --- Let us take for M the set \(R_{+}\) , equipped with the law of addition and with the usual topology. Let P be the set of complex numbers z with positive real part, equipped with the topology induced by C, and with the law induced by addition in C. For every \(p \in P\) , denote by \(f(p)\) the character \(x \mapsto e^{-px}\) of \(R_{+}\) ; it is easily verified that f is an isomorphism of the topological monoid structure of P onto that of X; we shall identify X with P by means of f.~It is clear that \(R_{+}\) is a full submonoid of P, and Th. 3 yields the following result:

PROPOSITION 16. --- Let A be a set equipped with a filter \(\mathfrak{F}\) . For every \(\alpha \in \mathbf{A}\) , let \(\mu_{\alpha}\) be a bounded positive measure on \(\mathbf{R}_{+}\) , and let \(\Phi_{\alpha}\) be the function defined on \(\mathbf{R}_{+}\) by \(\Phi_{\alpha}(p) = \int_{0}^{+\infty} e^{-px} d\mu_{\alpha}(x)\) . In order that the mapping \(\alpha \mapsto \mu_{\alpha}\) converge tightly with respect to \(\mathfrak{F}\) to a bounded positive measure \(\mu\) , it is necessary and sufficient that \(\Phi_{\alpha}\) converge pointwise on \(\mathbf{R}_{+}\) , with respect to \(\mathfrak{F}\) , to a function \(\Phi\) continuous at the point 0. In this case, \(\Phi(p) = \int_{0}^{+\infty} e^{-px} d\mu(x)\) for all \(p \in \mathbf{R}_{+}\) .

There are analogous results for the additive monoids \(\mathbf{R}_{+}^{n}\) (n an integer \textgreater1); we leave to the reader the transcription of Th. 3 in this case.

